\item The two real eigenspaces of $T\chi$ have rank one.
\item The theorem uses inverse Frobenius in its purity condition.
On $T\chi$ its good eigenvalues have complex absolute value
$q^{-1/2}$, so its integer $w$ is $-1$.
\item $V\chi$ is irreducible; twisting preserves irreducibility of
the non-CM rational Tate representation.
\item The image over $\mathbf Q(\mu_{2^\infty})$ contains an open
subgroup of $\mathrm{SL}_2(\mathbf Z_2)$. Intersecting with
$\ker\chi$ keeps an open subgroup, which contains a nonidentity
unipotent. It fixes a one-dimensional subspace and gives condition~(v).
\end{itemize}
These are all rational hypotheses. The residual irreducibility,
integral cokernel and $p\ne2$ conditions of part~(3) are not used.
Theorem~13.4(1) gives torsion $\mathbf H^2(T\chi)$, and part~(2)
gives, at every height-one prime away from two,
\[
 \ell\mathbf H^2(T\chi)_{0,\mathfrak p}
       \leq\ell(A_{\mathfrak p}/J_{\mathfrak p}),\qquad
 \mathbf H^2_0=\ker(\mathbf H^2\to\mathbf H^2_{\rm loc}). \tag{K15}
\]

The finiteness and Euler characteristic (12.2.1)--(12.2.2), p.~220,
give generic rank one in $\mathbf H^1$. Its localization away
from two is torsion-free. Indeed a torsion element for a scalar
prime would come from invariants of the corresponding finite-dimensional
rational coefficient algebra at that prime; its coefficient action is
abelian, whereas the open commutator subgroup on $V\chi$ has no
fixed vectors. The coefficient exact sequence excludes these
invariants. The scalar specialization need not be a finite-order
character. Thus $\mathbf H^1_{\mathfrak p}$ is free of rank one.
Localization commutes with $\operatorname{Hom}$ for this finitely
presented module. Writing $y_{2^\infty}=a e$ in a local basis gives
$J_{\mathfrak p}=(a)$. Formula (K15) is exactly the index bound for
the supplied class, not for an arbitrarily rescaled generator.
